% Part: first-order-logic
% Chapter: sequent-calculus
% Section: rules-and-proofs

\documentclass[../../../include/open-logic-section]{subfiles}

\begin{document}

\olfileid[pt]{mvl}{seq}{rul}

\olsection{Regras e \usetoken{P}{derivation}}

No que segue, sejam $\Gamma, \Delta, \Pi, \Lambda$ sequências finitas
de !!{sentence}s.

\begin{defn}[Sequente]
Um \emph{sequente de $n$ lados} é uma expressão da forma
\[
\Gamma_1 \nSequent \dots \nSequent \Gamma_n
\]
em que cada $\Gamma_1$ é uma sequência finita (possivelmente vazia) de
!!{sentence}s da linguagem~$\Lang L$.
\end{defn}

\begin{defn}[Sequente inicial]
Um \emph{sequente inicial de $n$ lados} é um sequente de $n$ lados da
forma $!A \nSequent \dots \nSequent !A$ para qualquer !!{sentence} $!A$
da linguagem.

Se a linguagem contém um conectivo de aridade~$0$, $\star$, isto é, uma
constante proposicional, tomamos também o sequente $\dots \nSequent
\star \nSequent \dots$, no qual $\star$ aparece na posição do valor de
verdade associado a~$\tf{\star} \in V$, e as demais posições ficam vazias.
\end{defn}

Para cada conectivo de uma lógica $n$-valorada~$\Log{L}$, há uma regra
lógica para cada valor de verdade que esse conectivo pode assumir
em~$\Log{L}$. As !!^{derivation}s em um cálculo de sequentes de $n$
lados para~$\Log{L}$ são árvores de sequentes, nas quais os sequentes
superiores são sequentes iniciais e, se um sequente aparece abaixo de
um ou mais sequentes, deve decorrer corretamente deles por uma regra
de inferência para os conectivos de~$\Log{L}$.

\begin{defn}[Teoremas]
!!^a{sentence}~$!A$ é um \emph{teorema} de uma lógica
$n$-valorada~$\Log{L}$ se há !!a{derivation} do $n$-sequente que contém
$!A$ em cada posição correspondente a um valor de verdade designado
de~$\Log{L}$. Escrevemos $\Proves[\Log{L}] !A$ se $!A$ é um teorema e
$\Proves/[\Log{L}] !A$ se não o é.
\end{defn}

\begin{defn}[!!^{derivability}]
!!^a{sentence}~$!A$ é \emph{!!{derivable} a partir de} um conjunto de
!!{sentence}s~$\Gamma$ em uma lógica $n$-valorada~$\Log{L}$, $\Gamma
\Proves[\Log{L}] !A$, sse há um subconjunto finito~$\Gamma_0 \subseteq
\Gamma$ e uma sequência $\Gamma_0'$ das !!{sentence}s em~$\Gamma_0$
tais que o sequente a seguir tem !!a{derivation}:
\[ \Lambda_1 \nSequent \dots \nSequent \Lambda_n \] em que $\Lambda_i$
é $!A$ se a posição $i$ corresponde a um valor de verdade designado, e
$\Gamma_0'$ caso contrário. Se $!A$ não é !!{derivable} a partir de
$\Gamma$, escrevemos $\Gamma \Proves/ !A$.
\end{defn}

Por exemplo, a lógica $3$-valorada de \L ukasiewicz tem um cálculo de
sequentes de $3$ lados. Em um sequente de $3$ lados $\Gamma \nSequent
\Pi \nSequent \Delta$, $\Gamma$ corresponde a~$\False$, $\Delta$
a~$\True$ e $\Pi$ a~$\Undef$. Os axiomas são $!A \nSequent !A
\nSequent !A$. Como apenas $\True$ é designado, $\Gamma
\Proves[\LogLuk[3]] !A$ sse o sequente $\Gamma \nSequent \Gamma
\nSequent !A$ tem !!a{derivation}. (Se $\Undef$ também fosse designado,
precisaríamos de !!a{derivation} de $\Gamma \nSequent !A \nSequent !A$.)

\end{document}
